\qhead{10. Price Leadership and Price Wars — Is there a price leader in your industry? How does it perform its role? Are the others followers or challengers? Do you see potential for a price war?}

\textbf{Is BYD a price leader?} Partly. Good pricing conduct needs \emph{visibility}, \emph{common motivation} and \emph{resolve} \parencite{nagle2018}. BYD has visibility (the largest, most-watched player) and resolve (it cuts credibly and sustains), but lacks \textbf{common motivation}: its cost edge means it \emph{wants} lower prices than rivals. So instead of leading the industry \emph{up}, it leads it \emph{down} — a price-war \textbf{instigator} (cuts above 30\% on some models in May 2025 \parencite{restofworldWar}), not a stabiliser.

\textbf{Potential for a war?} It is already underway, and the risk checklist is high on nearly every factor: stagnant growth, capacity glut, fixed-cost-heavy production, low differentiation, price-driven buyers, transparent prices, strong retaliation ability.

\textbf{Should a rival match?} For a higher-cost legacy OEM the response flowchart says: the cut is real and broad; large volume is at risk; but it \emph{cannot} sustain the war (higher break-even, Question~7). \textbf{Conclusion: do not match on price} — use non-price responses (reinforce differentiation, retreat to a defensible niche).

\textbf{Should BYD have started it?} The ``justified'' conditions fit — a sustainable cost advantage and the ability to outlast rivals — but even the apparent winner paid: profit fell 19\% and regulators stepped in \parencite{byd2025}.
